\documentclass[UTF8,fontset=none,11pt,a4paper]{ctexart}
\usepackage[left=21mm,right=21mm,top=27mm,bottom=22mm,headheight=26pt,headsep=6mm]{geometry}
\usepackage{fontspec}
\setmainfont{texgyretermes-regular.otf}[BoldFont=texgyretermes-bold.otf,ItalicFont=texgyretermes-italic.otf]
\setCJKmainfont{Songti.ttc}[Path=/System/Library/Fonts/Supplemental/,FontIndex=6,AutoFakeBold=2]
\setCJKsansfont{Songti.ttc}[Path=/System/Library/Fonts/Supplemental/,FontIndex=6,AutoFakeBold=2]
\usepackage[table]{xcolor}
\definecolor{TableHead}{HTML}{EEF0ED}
\definecolor{Muted}{HTML}{666C68}
\usepackage{booktabs,longtable,array,fancyhdr,needspace,xurl}
\usepackage[unicode,hidelinks]{hyperref}
\urlstyle{same}
\linespread{1.32}
\setlength{\parindent}{2em}
\setlength{\parskip}{3pt}
\setlength{\LTpre}{8pt}\setlength{\LTpost}{10pt}
\ctexset{section={format=\fontsize{15}{20}\selectfont\bfseries,beforeskip=18pt,afterskip=8pt},subsection={format=\fontsize{12}{16}\selectfont\bfseries,beforeskip=12pt,afterskip=6pt},subsubsection={format=\normalsize\bfseries,beforeskip=9pt,afterskip=4pt}}
\pagestyle{fancy}\fancyhf{}
\renewcommand{\headrulewidth}{.3pt}
\renewcommand{\footrulewidth}{0pt}
\fancyhead[L]{\parbox[b]{\linewidth}{\fontsize{8}{10}\selectfont\color{Muted}投资账户：自动管理还是自行管理·自动投资项目披露}}
\fancyfoot[C]{\small\color{Muted}\thepage}
\hypersetup{pdftitle={自动投资项目披露},pdfsubject={仅供模拟使用·投资账户：自动管理还是自行管理·16-01},pdfauthor={金融决策模拟材料}}
\begin{document}
{\noindent\fontsize{20}{25}\selectfont\bfseries 自动投资项目披露\par}
\vspace{7pt}{\noindent\small\color{Muted}仅供模拟使用·投资账户：自动管理还是自行管理·16-01\par}\vspace{10pt}
本文件为自动投资项目披露资料，说明项目资格及业务做法。资料日期2026年6月30日。


\section*{引言}
经纪公司兼投资顾问提供自动投资项目，以及其他综合管理项目。自动项目按客户陈述的投资目标与风险承受能力，提供多元化组合及全权管理。


\section*{服务、费用与报酬}
项目通过自动投资顾问服务提供全权管理。组合包括ETF及按目标和风险承受能力决定的现金配置。现金存放于关联银行的存款账户。FDIC保障须满足条件，经纪公司及投资管理公司不是银行；非存款产品不受FDIC保障，可能损失价值。


ETF可以涵盖股票、固定收益、房地产及商品。资产类别数量及类型可能随时调整，不事先通知。项目设计为每日监测，按需要自动再平衡，使组合与选择的风险档案一致，除非认为再平衡不符合客户利益。


项目通过网站及手机应用提供。客户回答在线问卷，系统判断适用性并给出建议组合，供客户审核批准。加入前，应考虑项目及组合是否符合需要与目标；加入后可以更新投资档案，查看配置、活动及表现。目标、风险承受能力或财务状况变化时，应定期检查和更新。


关联投资管理公司按客户选择策略提供全权组合管理。经纪公司是合格托管人，提供交易执行、行政及相关服务。


算法是电脑程序中的规则。项目算法建议组合、识别再平衡机会、识别税损收割机会，以及发起有关买卖指令。再平衡算法设计为每日检查账户及持仓，指令先送投资管理公司审核，再交付执行。技术及其他情况可能使算法不能运行或不能执行原定活动。


算法包含洗售避免功能，在一些情况暂停再平衡及税损收割，可能超过三十日不能回到目标配置，并错失机会。


客户授予管理公司交易及管理权。公司可以自行发起或暂停交易，包括认为继续交易可能不当损害账户时。客户不能自行在项目账户交易。


可请求排除某些ETF，但不一定全部处理。股票、共同基金或现金等其他限制如被认为不合理，不予执行；ETF限制可能延迟管理。无法按限制管理时，会通知客户。


税损收割需请求并满足单一账户目前50,000美元门槛，该门槛可能改变。一般不宜借用融资余额加入项目。项目账户不是融资账户，不能作融资贷款抵押。


\subsection*{经纪做法}
本项目不由经纪公司作为顾问另选经纪商。客户同意全部交易由本经纪公司执行，因此不一定取得其他经纪商可能提供的最有利价格或执行。并非所有经纪商或顾问都要求指定交易机构。


\subsection*{托管}
公司保管项目资产，至少每季度提供报表，列明持仓与期间活动。客户应认真检查。


\subsection*{全权投资及电子交付}
\textbf{客户取得全权ETF与现金组合管理。}加入时同意电子签署，以及电子收取记录、披露、补充文件和账户文件。


客户须维持准确有效邮箱，并能够阅读、下载、打印及保存电子文件。不愿或不能电子接收时，公司可终止项目，可能要求转移到其他本公司或其他机构账户，或清算资产退回款项。


\subsection*{项目费用}
\textbf{服务包含行政、交易执行、托管及管理。}客户不直接支付年度项目费，但项目不是没有成本。客户承担ETF运营费，表现也受现金配置及归集影响。公司及关联机构从部分ETF及现金安排取得报酬。


类似服务中，客户可能预期支付资产0.30\%的年费。由于项目含退休账户，公司按IRS规则进行名义计算，以关联机构从ETF交易获得的报酬抵销0.30\%；超过资产0.30\%的有关报酬，会退还或用于账户行政成本。结果是不收年度项目费。


\subsection*{其他费用}
激活或使用税损收割不另收规划或顾问费，项目交易不收经纪佣金。ETF支付顾问、行政、分销、转移代理、托管、法律、审计等费用，最终由股东承担，关联管理公司从关联ETF取得管理费。


出售原持仓或办理其他证券交易，可能产生销售、赎回费用及税务后果。管理公司不与项目客户直接签合同或协商费用。


银行向经纪公司支付每账户固定行政服务年费。组合服务成本可能高于或低于分别购买；比较因素包括账户规模、可合并账户、交易数量、投资及现金种类，以及额外顾问服务范围。


\subsection*{报酬}
分行及电话财务顾问通常推荐项目。顾问可能是雇员，也可能是经营独立分行的承包商。雇员在基本薪资之外，可因引导客户加入顾问项目及后续服务而取得报酬。


报酬基于理解、推荐及服务所需时间、复杂度及专业知识。公司因客户选择不同项目获得不同收入；顾问在项目中资产取得的报酬可高于佣金经纪账户。独立分行负责人从公司获得月净付款，再向其雇员支付薪酬，不同项目金额不同。


其他投资人员也可能因转介顾问服务取得额外激励。吸引净新增资产还可能获得年度旅行奖励，不论资产是否加入自动项目。


\subsection*{关联机构利益及利益冲突}
收入包括银行从现金配置获得的收入，管理公司从选用的关联ETF取得的顾问费，以及交易订单路由市场中心可能支付的报酬。


项目不收顾问费，部分依靠现金配置带来的银行收入。作为项目提供者，经纪公司设定各策略现金比例，并考虑包括银行预期收入在内的因素，形成利益冲突。其他风险资产表现优于现金时，现金配置可能降低总体表现。


现金比例越高、向客户支付的利率越低，银行收入越多。较低利率降低现金收益。这是利益冲突，项目书面参数并不能消除它。项目以外可取得收益较高的现金产品，但可能没有FDIC保障。


银行净利息收入来自贷款及投资证券取得的利息，减向归集客户支付的利息，差额为利差。市场利率上升，历史上利差通常增加。现金比例高于10\%时，银行通常收入较多；低于10\%时较少。


经纪公司和管理公司关联，存在选择和保留关联管理公司的激励。管理公司选择向其付费的关联ETF也存在冲突。经纪公司指示在满足ETF标准时选用或保留关联ETF。


部分资产类别同时使用市值加权和基本面加权ETF。前者按市值，后者按销售、现金流、股息和回购等因素加权，以促进分散和调整风险收益。基本面ETF通常费率较高；目前的第三方分类与选择方法，使关联ETF成为基本面类别的主要选择。


除ETF选择部分所述情形外，书面参数不允许管理公司考虑向经纪公司或其他关联方支付的报酬。管理公司须按书面标准选择和移除证券，经纪公司检视其服务表现。加入前，应考虑这些冲突对管理的影响。


\section*{账户要求与客户类型}
项目面向希望自动管理、接受主要通过数字渠道服务的投资者。可包括个人、IRA、生前信托及某些组织账户；政府实体及受ERISA规管的客户没有资格。不同组织可能有不同客户体验。


客户须同意电子收取合同、披露、说明书、报表等。每账户至少投入5,000美元才能开始投资。税损收割另有启动及维持余额要求。


可用证券入金。混合符合及不符合建议组合的证券，可能产生额外再平衡交易。客户授权公司出售投入的证券，不考虑因此实现的应税损益；清算前表现不由公司负责。


客户可按顾问协议随时退出。退出后现金归集项目不再适用，账户改用自己的归集安排，其条件可能较好或较差。


公司可以因未达到最低资金、邮箱无效、不合理投资限制或其他原因终止客户。依原因可能允许补救。任一方终止后，归集终止，账户不再受项目管理。


\subsection*{退休账户及管理公司选择}
经纪公司不对退休项目资产行使投资全权，也不自行作全权行为，仅执行有关指令；投资管理公司是执行退休账户管理及全权投资的唯一受托人，受合理客户限制约束。


公司选择关联管理公司，并通过季度标准组合表现报告评估。关联公司也管理其他策略及基金，管理采用相同证券的不同策略可能有冲突。


公司不因转介本项目取得按表现收费。不由经纪公司自己管理或推荐策略，因此其并行管理项目不适用。


\section*{分析方法、策略及损失风险}
管理公司按经纪公司配置和ETF参数形成多种策略。策略组合ETF与现金归集；有些用于应税账户，有些用于递延税账户，有些侧重利息及股息收益类别。策略或参数可能不预告地增加或修改，并非所有策略适用所有客户。


\textbf{投资证券有损失风险。}算法在既定参数及限制内维持分散，不按短期市场波动主动调整配置，也不考虑全部短期波动或个人资本利得税情况。人员监督算法，不亲自逐一监测各账户。


算法或软件可能不按预期参数运行，导致建议组合过于积极或保守，或错误触发／漏发再平衡与税损交易。算法互相作用并按优先顺序执行；计划提款优先于避免洗售，必要出售即使造成洗售也会执行。


\subsection*{网络安全、人工智能及系统中断}
互联网、手机应用及云服务可能遭受病毒、勒索软件、钓鱼或过载攻击，导致个人财务资料未经授权访问、数据损失修改、账户被接管或暂时无法使用。安全措施和外部专家不能保证阻止所有不断改变的攻击。


公司使用机器学习及自然语言等AI技术，可能依赖内部或第三方数据模型；部分系统不完全受直接控制，可能提供错误、不完整、有偏或暂时不可用的资料。涉及网络安全、隐私、可靠性及运营中断风险。快速变化的法律可能要求修改使用方式并增加成本。控制、审核及人工监督不能消除全部风险。


软件错误、停电、使用量、第三方、天气、灾害或攻击可能使系统中断，限制查看交易，影响管理或服务，并需要时间修复。


\subsection*{ETF选择}
符合资格的ETF应代表组合资产类别，具足够流动性，并属于其类别中低运营费产品。替换时还考虑客户新增交易及其他成本。满足因素时，公司指示选择或保留关联ETF。


通常每类选主要和次要ETF，考虑税损收割及限制；某些情况仅一种ETF，此类别不能执行税损收割。主要ETF股价还需低于适合小账户交易的上限。


\subsection*{再平衡}
每日检查偏离目标是否超过参数，符合时发起指令，除非管理公司因市场波动、税损或洗售等决定不执行。资料变更、投资限制及存取款也可能触发。通过ETF买卖及归集存取恢复配置。技术限制可能使监测或交易不在每日发生。


余额低于5,000美元，可能更偏离参数及目标配置；因无法买卖至少一份ETF，提款可能需要卖出整个资产类别。


\subsection*{洗售避免}
卖亏损证券前后三十日买入相同或实质相同证券，可能使IRS不认可或递延亏损扣除。算法可能暂停交易三十日以上，影响配置或错失税损机会。


对同一主要账户持有人关联的项目账户，除部分托管、信托及未成年继承IRA等例外，尝试共同避免洗售。首个应税账户亏损出售某ETF，会阻止其他账户三十日内买入；首先买入，则阻止其他应税账户三十日内亏损出售。


算法不能避免全部洗售。提款所需出售仍执行；用亏损卖出款入金后，项目三十日内买入同ETF也可能造成洗售。只监测项目账户，客户须检查自己及配偶在项目之外的账户。


\subsection*{税损收割功能}
启动须应税账户达到50,000美元，维持另有较低门槛。低于维持门槛时不活跃，直到再次符合。门槛降低但不能完全避免替代ETF单份价格过高，出售款部分留在现金的情况。


达到门槛时，算法按日检查机会并发起交易，但受资产类别和洗售限制。亏损出售以抵销潜在资本收益；公司不监测资本收益种类及金额。类别比重不足时，选用依不同指数判断实质不同的替代ETF。


替代持仓有亏损批次时，通常三十日不能出售，即使其他批次有盈利。因此再平衡可能不能按策略执行。多个共同监测账户在同类持有不同ETF时，优先在最早加入账户执行，跳过其他账户，不优先选择产生最大亏损的ETF。


替代ETF表现可能优于或低于被出售ETF。亏损的使用取决于同年或未来资本收益及税法限制。效果受完整税务投资状况，包括配偶或外部账户、投资种类及持有期影响。不保证降低、递延或消除任何年度税负。应向专业税务顾问咨询；公司不保证特定税务结果。


\subsection*{自动投资现金归集项目}
各策略现金比例通常在账户6\%至30\%之间，依目标、风险与期限决定。加入时同意未投资贷方余额归集至关联银行存款。


归集是必需功能。现金超过配置及再平衡参数时，用超额买入证券。提款可能需要出售ETF，以恢复现金配置；应税账户可能产生实现损益。


客户授权公司将现金安排改为其自身、关联方或第三方现金储蓄或投资工具；具体条件见归集披露。


自2022年6月1日起，银行同意按月将现金利率设为指定政府货币基金归集份额SWGXX的七日收益率（包含减免）。利率及计算方法均可改变。银行按协议或监管要求通知经纪公司改变方法，公司收到后通知客户。


该利率可能高于或低于同银行、其他银行存款，以及经纪平台上货币基金等现金等价产品。银行以现金支持现有及新贷款和投资，其收入取决于贷款投资收益与付给客户利息的利差；资金也提供流动性、稳定及低成本融资。


\subsection*{等待投资或分配的资金}
入金或其他未投资余额在归集前，公司可能取得利息。营业日截止前收到，通常至下一营业日开始；非营业日或截止后收到，通常至第二个随后营业日开始。


提款支票通常处理两三个营业日，死亡或离婚等特殊情况可能更久。处理完成当天扣账，下一营业日写支票并邮寄。公司自扣账至支票被提出付款期间取得利息；客户可询问未兑现支票状态。


电子转账开始当天扣账，其他金融机构通常一至两个营业日收到；公司从扣账至收款机构收到、资金自公司银行账户划出期间取得利息。


\section*{投票、客户资料及联系}
经纪公司不接受证券代理投票权，管理公司协调项目投票。加入时及以后按需要，向管理公司提供策略及合理限制资料。


希望联系管理公司，可向经纪公司代表提出请求。经纪公司及代表为项目主要联系人。


\section*{其他金融行业活动及关联方}
经纪公司兼具经纪商及登记投资顾问身份。经纪业务提供证券信息、讲座、期刊、报告、指南、规划工具、保险资料及线上数据。其他顾问项目采用不同管理、转介及规划服务；参加其他项目的费用和条件不因本自动项目而自动适用。


母公司全资子公司经营投资顾问、经纪、信托、托管及银行业务。关联管理公司还管理其他全权策略、基金及ETF。非全权财富顾问提供定期建议、由客户作投资决定；其他子公司向合资格客户提供国际证券、退休计划、信托、期货、外汇、存款及贷款服务。


\section*{道德守则}
公司按投资顾问法SEC Rule 204A-1制定守则，反映顾问的受托原则。涵盖人员遵守证券法、报告违反守则的情形；有投资建议或相关资料访问权限者，定期报告某些个人交易及持仓。


访问人员不得向外部披露组合交易或非公开资料，执行客户交易所需除外；不得为自己或他人谋利，也不得在客户证券买卖中进行欺骗。守则可随监管或内部政策更新，可申请副本。


\subsection*{订单路由与执行}
公司向无关联经纪商及主要交易所发送非指定交易订单。评估价格与基准、价格改善、市场深度、订单规模、速度及准确度、处理系统、流动性、自动执行保障、限价可成交时执行概率、服务及成本。


价格改善是成交价优于全国最佳买卖报价。公司持续检查市场长期执行质量。项目证券通常不由公司作交易对手，但某些零碎份额例外。


零碎份额不能直接在市场出售或通过自动清算转移，公司可作为交易对手买入。这可能不是市场最佳价格，而且银行与客户处于交易两方，形成冲突；不另收零碎交易费用。


路由市场可能支付流动性或订单流返还，也可能收取执行费。季度市场及支付记录在网站提供，可申请书面资料。本人订单过去最多六个月的执行时间及去向，也可申请。


\subsection*{ETF及有关付款}
每只ETF依自身说明书支付投资顾问、行政、分销、转移代理、托管、法律、审计等费用，由股东承担；关联管理公司从项目所持关联ETF取得管理费。


第三方ETF提供者或关联方可能为教育、活动、数据及报告向公司付款。某些提供者每年为数据报告支付最高4,000,000美元；教育活动费用依种类和次数不同。关联管理公司也可能支付教育支持、资料报告及其他服务报酬。


\subsection*{员工个人交易}
公司监测代表证券持仓及交易，检查本公司及适用外部账户，按道德守则、合规手册及政策监督。雇佣后须披露拥有或控制的证券账户，定期检查确认。


禁止违反联邦、州及交易所规则，包括抢先客户交易、将保密客户交易信息传给他人谋利，以及购买首次公开发行股份。


\section*{账户审查}
加入项目时开立经纪账户，收到交易电子确认及季度报表。报表显示全部项目账户持仓活动、证券买卖、利息股息及其他收入、入提款、费用和账户价值，部分资料可能来自第三方。


你负责及时检查确认和报表，管理有疑问时直接联系公司。公司至少每年通过电子渠道请求在重大变化时更新信息。目标、财务、投资意图或合理限制改变，应及时在网站更新。未更新档案，公司不会自行改变所选组合。


\section*{客户转介与其他报酬}
雇员、独立分行及代表按费用报酬条款取得收入。另有USAA转介协议，按转介至本公司账户的资产，包括本项目，支付年度0.15\%，即十五个基点的报酬。


\section*{财务信息}
公司不要求或招揽预付费用，因此不必在此附上最近年度资产负债表。公司表示没有合理可能妨碍合同义务的财务状况，也没有当前或过去十年的破产申请。

\end{document}